\section{De la estructura compleja interna a la geometría proyectiva de Bloch}
\label{sec:geometria-proyectiva-bloch-rev2}
\index{esfera de Bloch}
\index{estructura compleja!realificación ortogonal}
\index{grupo icosaédrico binario}
\index{Möbius!acción proyectiva}

La raíz interna de menos la identidad posee dos realizaciones coordinadas. En
la Parte I, la incidencia de Paley--Witt construye
\(J_W^2=-I\) sobre \(\mathbb F_3^6\), organiza el código ternario y realiza
la estructura \(\mathbb F_9^3\); en una realización analítica, un operador
ortogonal \(J_E\) con \(J_E^2=-I\) convierte un espacio real de dimensión
par en un espacio complejo. Esta segunda flecha añade topología, norma y
completación: conserva la relación cuadrática, pero no identifica los cuerpos
de escalares.

El tránsito completo queda tipado por
\begin{equation}
 \boxed{
 \begin{gathered}
 C_P\longrightarrow J_W^2=-I_6
 \longrightarrow \mathbb F_3^6\simeq\mathbb F_9^3,\\
 (E_4,J_E),\quad \dim_{\mathbb R}E_4=4,
 \longrightarrow
 \mathbb P(\mathcal H_J)\simeq\mathbb {CP}^{1}
 \longrightarrow S^2,\\
 \mathrm{SU}(2)\longrightarrow\mathrm{PSU}(2)\simeq\mathrm{SO}(3)
 \subset\mathrm{PSL}(2,\mathbb C)\cong\text{\rmfamily Möb}^{+}(S^2),\\
 2A_5=\pi^{-1}(A_5)\subset\mathrm{SU}(2).
 \end{gathered}}
 \label{eq:cadena-paley-bloch-rev2}
\end{equation}
La primera línea conserva sus pruebas en los
\cref{chap:paley-codigo-witt,thm:f9-tres-desde-paley}; aquí se construyen las
dos líneas proyectivas.

\subsection{Realificación de la fase y espacio de rayos}

Sea \((E,g)\) un espacio euclídeo real de dimensión par y sea
\(J_E\in\operatorname{End}(E)\) tal que
\[
 J_E^2=-I,
 \qquad g(J_Eu,J_Ev)=g(u,v).
\]
Entonces \(J_E^{\mathsf T}=-J_E\) y
\[
 R_J(\theta):=e^{\theta J_E}
 =\cos\theta\,I+\sin\theta\,J_E
\]
es ortogonal, con
\(R_J(\theta)R_J(\phi)=R_J(\theta+\phi)\). La acción
\((a+ib)u:=au+bJ_Eu\) realiza la carta compleja sin utilizar el escalar
\(i\) como antecedente del operador.

Sea ahora \(E_4\subseteq E\) un subespacio \(J_E\)-estable de dimensión real
cuatro y sea \(\mathcal H_J=(E_4,J_E)\) el espacio de Hilbert resultante, de
dimensión compleja dos.
Un vector unitario
\[
 \psi=\begin{pmatrix}z_0\\z_1\end{pmatrix},
 \qquad |z_0|^2+|z_1|^2=1,
\]
determina un rayo \([\psi]\in\mathbb {CP}^{1}\). Definamos
\begin{equation}
 r_1=2\operatorname{Re}(\overline z_0z_1),
 \qquad
 r_2=2\operatorname{Im}(\overline z_0z_1),
 \qquad
 r_3=|z_0|^2-|z_1|^2.
 \label{eq:coordenadas-bloch-rev2}
\end{equation}

\begin{theorem}[Identificación proyectiva de Bloch]
\label{thm:bloch-proyectivo-rev2}
La aplicación
\[
 \mathcal B:\mathbb {CP}^{1}\longrightarrow S^2,
 \qquad [\psi]\longmapsto(r_1,r_2,r_3),
\]
es una biyección suave. El proyector de rango uno asociado al rayo es
\begin{equation}
 \rho_\psi=|\psi\rangle\langle\psi|
 =\frac12\bigl(I_2+r_1\sigma_1+r_2\sigma_2+r_3\sigma_3\bigr),
 \label{eq:proyector-bloch-rev2}
\end{equation}
donde \(\sigma_1,\sigma_2,\sigma_3\) son las matrices de Pauli.
\end{theorem}

\begin{proof}
La fase global deja invariantes las tres coordenadas y
\[
 r_1^2+r_2^2+r_3^2
 =4|z_0|^2|z_1|^2+(|z_0|^2-|z_1|^2)^2=1.
\]
Recíprocamente, cada \(\boldsymbol r\in S^2\) determina el proyector positivo
\(\rho=(I+\boldsymbol r\cdot\boldsymbol\sigma)/2\), que satisface
\(\rho^2=\rho\) y \(\operatorname{tr}\rho=1\); su imagen es un único rayo.
Las cartas afines de \(\mathbb {CP}^{1}\) hacen suaves el mapa y su inversa
\cite{Bloch1946}.
\end{proof}

La secuencia de Hopf
\begin{equation}
 S^1\longrightarrow S^3\longrightarrow S^2
 \label{eq:fibrado-hopf-bloch-rev2}
\end{equation}
expone la pérdida publicada por la proyectivización: el vector normalizado
pertenece a \(S^3\), el observable proyectivo publica un punto de \(S^2\) y
la fibra \(S^1\) conserva la fase común que la carta identifica
\cite{Hopf1931}. Esta es una realización precisa de la diferencia HMT entre
estado enriquecido y coordenada visible.

\subsection{Rotaciones, transformaciones de Möbius y elevación espinorial}

Todo elemento de \(\mathrm{SU}(2)\) puede escribirse
\[
 U=\begin{pmatrix}
 \alpha&\beta\\-\overline\beta&\overline\alpha
 \end{pmatrix},
 \qquad |\alpha|^2+|\beta|^2=1.
\]

\begin{theorem}[Acción espinorial sobre la esfera de Bloch]
\label{thm:su2-so3-bloch-rev2}
Para cada \(U\in\mathrm{SU}(2)\) existe una única rotación
\(R_U\in\mathrm{SO}(3)\) tal que
\[
 U(\boldsymbol r\cdot\boldsymbol\sigma)U^\dagger
 =(R_U\boldsymbol r)\cdot\boldsymbol\sigma.
\]
La aplicación \(U\mapsto R_U\) es un homomorfismo sobreyectivo con núcleo
\(\{\pm I_2\}\).
\end{theorem}

\begin{proof}
La conjugación conserva el espacio real tridimensional de matrices
hermíticas sin traza y el producto escalar
\(\langle A,B\rangle=\tfrac12\operatorname{tr}(AB)\), por lo que induce una
rotación. Si actúa trivialmente, \(U\) conmuta con las tres matrices de Pauli
y es escalar; el determinante uno deja \(U=\pm I_2\). La representación
eje--ángulo proporciona las dos preimágenes de cada rotación.
\end{proof}

En la carta afín \(z=z_0/z_1\) de \(\mathbb {CP}^{1}\), la misma acción es
la transformación de Möbius
\[
 z\longmapsto
 \frac{\alpha z+\beta}{-\overline\beta z+\overline\alpha}.
\]
La rotación de Bloch y la transformación fraccional son dos coordenadas de
una única acción proyectiva: la primera actúa sobre proyectores hermíticos y
la segunda sobre la esfera de Riemann.

El grupo rotacional del icosaedro es \(A_5\subset\mathrm{SO}(3)\). Su
preimagen define el grupo icosaédrico binario
\[
 2A_5:=\pi^{-1}(A_5)\subset\mathrm{SU}(2),
\]
y produce la sucesión exacta
\begin{equation}
 1\longrightarrow\{\pm I_2\}
 \longrightarrow2A_5\longrightarrow A_5\longrightarrow1.
 \label{eq:extension-central-2a5-rev2}
\end{equation}
En particular, \(|2A_5|=120\). Para una rotación de orden impar \(n\), el
levantamiento de semiángulo posee orden \(2n\), mientras su opuesto posee
orden \(n\). Las rotaciones icosaédricas de órdenes tres y cinco admiten, por
tanto, pares de levantamientos de órdenes \((6,3)\) y \((10,5)\),
respectivamente. Esta especialización enlaza la
geometría proyectiva de los estados de dos niveles con la rama de McKay y las
simetrías excepcionales, cuyos módulos y retículas permanecen en sus
construcciones de referencia de la Parte I.

\begin{statusbox}[title={Procedencia y alcance de la cadena proyectiva}]
La identidad de Paley--Witt, la estructura
\(\mathbb F_3^6\simeq\mathbb F_9^3\), la acción
\(\mathrm{SU}(2)\to\mathrm{SO}(3)\), la realización de Möbius y la extensión
\(2A_5\) son resultados demostrados en sus dominios. La cadena
\eqref{eq:cadena-paley-bloch-rev2} es una realización functorial unificada:
declara explícitamente la realización que cambia de escalares y no convierte
la raíz finita \(J_W\) en un operador analítico sin ese dato adicional.
\end{statusbox}
